

\section{Eleventh Sunday after Trinity: Self-Exaltation and Self-Humbling}


\begin{quote}

Matthew 23:1–12. Then Jesus spoke to the people and to His disciples, saying: The scribes and the Pharisees sit on Moses’ seat. Therefore, whatever they tell you to observe, observe and do; but do not according to their works. For they say it well enough, but they do not do it. For they bind heavy burdens, such as are hard to bear, and lay them upon men’s shoulders; but they themselves will not lift a finger to move them. But all their works they do in order to be seen by men; for they make their phylacteries broad and the tassels on their garments large. And they love the places of honor at feasts and in the chief seats in the synagogues. And they love to be greeted in the marketplaces and to be called by men, Rabbi, Rabbi. But do not let yourselves be called Rabbi; for one is your master, Christ, but you are all brothers. And you shall call no man on earth your Father; for one is your Father, He who is in Heaven. And neither should you allow yourselves to be called teacher; for one is your guide, Christ. But the greatest among you shall be your servant. But whoever exalts himself shall be humbled, and whoever humbles himself shall be exalted.

\end{quote}

\bigskip


“I thank You, God, that I am not as other men are,” says the Pharisee who went up to the temple to pray. This saying lays bare, in a single stroke, the whole Pharisaic mind. It appears so right and proper in the eyes of men that one tries to rise above another, that we are almost inclined to regard it as a virtue, contrary to all law and order, if anyone dares, together with our Lord and Savior, to speak the naked truth: “What is high among men is an abomination before God.”

Can it really be true? Has not an honorable man the right to praise himself for his honor? Has not an upright man reason to praise himself for his uprightness? Has not a rich man reason to be proud of the wealth he has acquired by diligence? Is it not altogether fair to despise thieves, drunkards, robbers, deceivers, adulterers, and ruffians? Thus one may ask, almost without end; and the answer comes so naturally and so straightforwardly from the mouth of the world that it seems almost impossible to object to it.

But the Savior nevertheless dares to do so. He says: “You are those who justify yourselves before men, but God knows your hearts.” He will not endure the false appearance that conceals the world’s proud thoughts. Among men such pride is well regarded, if only it keeps within reasonable bounds; but before God it is nonetheless an abomination.

But within God’s Church, much of that same spirit returns again, to our sorrow. There are many who have fallen into the Pharisees’ and the scribes’ dreadful sin. Their godliness has become for them a matter of pride and contempt. They will not count themselves as one with the world, and as they forget love, they become, by the new pride with which they despise the world, once again equal with the world in their hearts.

And those who sit on Moses’ seat, who are teachers and leaders in God’s congregation, how often do they not succumb to the terrible temptation which drove the Pharisees and the scribes into pride and self-exaltation. The honor that belongs to God and His Word they steal for themselves and seize to themselves. And the congregations are only too often willing to have it so. There are many members of the congregation who would rather hold to the priest than bow before the Word; and there are many members of the congregation who would rather go out of the congregation when they have fallen out with the priest than follow Jesus’ command and bidding, to heed the Word and treat the person as a separate matter.

From this there has come an abundance of misery also among us. A priest grows great in his own eyes because his calling is high and holy, and he wants to rule over the Lord’s flock instead of serving it with the Lord’s Word and the Lord’s love. By his lust to rule and his grand self-importance he provokes resistance and strife. Then the priest’s opponents are to be God’s opponents, and thus the congregation is torn apart. And conversely: members of the congregation come into conflict with the priest, and then they use their personal ill will toward the priest as an excuse for their own impenitence and stubborn resistance to the Word, and then either the priest must go, or else they must go out of the congregation.

Or can anyone deny that this is often how it happens?

O, how difficult it is to be a true priest and a true member of the congregation! To esteem God’s Word highly and oneself lowly, to set the Lord high and oneself low, when shall we ever learn it rightly?

There is only one way in which it can be learned. It is a very narrow way; for it is the Holy Spirit’s work. If any other spirit reigns in priest and congregation than the Holy Spirit, then it is the spirit of the world. But the spirit of the world always leads to self-exaltation and pride. If God’s congregation is built without God’s Spirit, it is gathered either by the priest’s brilliant gifts or by the priest’s cleverness and cunning or by the congregation members’ worldly competence, and the Spirit from God is lacking—then human self-exaltation and pride must follow, however much a show of godliness may be cast over it all.

If, therefore, when we build God’s Church on earth, we are not to fall back into the Pharisees’ way and the darkness of the papacy, then we must heed Jesus’ word concerning the scribes and the Pharisees. It concerns the priests, and it concerns the congregations. There is only one way for both to right and true humility before God and to true, living love toward one another. It is, as said above, the Holy Spirit’s way. There is only one Lord, that is Jesus Christ, but “no one can call Jesus Lord except by the Holy Spirit.” There is only one Father, namely He who is in Heaven, but no one can call Him Father except by the Spirit of sonship. There is only one teacher, that is Christ, and no one can be Christ’s follower through the cross to glory except the one who is driven by Christ’s Spirit. Therefore you are all brothers, and no one has a higher rank or is of nobler birth than another. In this is the Christians’ equality; in this is their brotherhood. This is the realization that must penetrate all God’s children: By grace, by mercy, by free unmerited love, God has saved me—saved us all; therefore I have nothing of which to boast.

But here lies Christian freedom, that the greatest is willing to be the servant of all. Not that the priest shall command the congregation, nor that the congregation shall command the priest; but that all should be willing to serve one another, and the strongest and greatest shall be the most eager to take up the heaviest burdens, the most willing to bear the heaviest blow, the most steadfast in labor, the most patient in suffering.

But flesh and blood cannot grasp that kind of greatness and that kind of freedom. It appears a contradiction to our reason that greatness and freedom should consist in humility and self-humbling. Then we must look to Jesus Christ, who left us an example, that we should follow His footsteps.

He understood it: both to speak the truth without partiality and to suffer without murmuring whatever fearless witness to the truth might bring. He did not bend the Lord’s truth in order himself to gain exaltation and honor; but He did not grow bitter against those who for that very reason nailed Him to the tree of the cross. He did not count equality with God a prize to be grasped, but willingly became the Lamb of God, who bore the sin of the whole world, and humbled Himself, so that He became obedient unto death, indeed the death of the cross. He sought not His own honor, but the Father’s; He sought not His own good, but the Church’s.

O, what an example to follow! Therefore let each one strike his breast and say: “God be merciful to me, a sinner!” For indeed, there is no one who has wholly followed the example; there is no one who has not shrunk back before the thorns on the narrow way, there is no one who does not carry within him a proud flesh, eager to exalt itself and despise others; there is no one who is not reluctant to be the servant of all.

But the Lord’s word stands fast forever: “Whoever exalts himself shall be humbled, and whoever humbles himself shall be exalted.”

